\documentclass[aspectratio=169]{beamer}
\usetheme{Madrid}
\usecolortheme{default}
\usepackage{amsmath,amssymb,bm}
\usepackage{braket}
\usepackage{booktabs}
\title{Variational Quantum Eigensolver with UCCSD}
\subtitle{Algorithm, Equations, and Workflow}
\author{MHJ, UiO}
\date{\today}
\begin{document}
%------------------------------------------------
\begin{frame}
\titlepage
\end{frame}
%------------------------------------------------
\begin{frame}{Goal}
We want to approximate the ground-state energy of an electronic Hamiltonian
\[
E_0 = \min_{\ket{\Psi}} 
\frac{\bra{\Psi} H \ket{\Psi}}{\braket{\Psi|\Psi}}.
\]
\vspace{0.3cm}
In VQE, we restrict the wave function to a parameterized ansatz
\[
\ket{\Psi(\bm{\theta})}
=
U(\bm{\theta})\ket{\Phi_0},
\]
and minimize
\[
E(\bm{\theta})
=
\bra{\Psi(\bm{\theta})}H\ket{\Psi(\bm{\theta})}.
\]
\vspace{0.3cm}
For UCCSD, the reference state is usually the Hartree--Fock state
\[
\ket{\Phi_0} = \ket{\Phi_{\mathrm{HF}}}.
\]
\end{frame}
%------------------------------------------------
\begin{frame}{Electronic Hamiltonian}
In second quantization,
\[
H
=
\sum_{pq} h_{pq} a_p^\dagger a_q
+
\frac{1}{2}
\sum_{pqrs}
h_{pqrs}
a_p^\dagger a_q^\dagger a_s a_r .
\]
\vspace{0.3cm}
Here
\[
h_{pq}
=
\int d\mathbf{x}\,
\phi_p^*(\mathbf{x})
\left(
-\frac12\nabla^2
-
\sum_A \frac{Z_A}{|\mathbf{r}-\mathbf{R}_A|}
\right)
\phi_q(\mathbf{x}),
\]
and
\[
h_{pqrs}
=
\int d\mathbf{x}_1 d\mathbf{x}_2\,
\frac{
\phi_p^*(\mathbf{x}_1)
\phi_q^*(\mathbf{x}_2)
\phi_r(\mathbf{x}_1)
\phi_s(\mathbf{x}_2)}
{|\mathbf{r}_1-\mathbf{r}_2|}.
\]
\end{frame}
%------------------------------------------------
\begin{frame}{Mapping Fermions to Qubits}
The fermionic Hamiltonian is mapped to qubits:
\[
H
\longrightarrow
H_q
=
\sum_{\alpha} c_\alpha P_\alpha ,
\]
where each
\[
P_\alpha
=
\sigma_1^{(\alpha)}
\otimes
\sigma_2^{(\alpha)}
\otimes
\cdots
\otimes
\sigma_n^{(\alpha)}
\]
is a Pauli string.
\vspace{0.3cm}
Common mappings include:
\begin{itemize}
\item Jordan--Wigner transformation,
\item Bravyi--Kitaev transformation,
\item parity mapping.
\end{itemize}
\vspace{0.3cm}
The VQE objective becomes
\[
E(\bm{\theta})
=
\sum_\alpha c_\alpha
\langle P_\alpha\rangle_{\bm{\theta}}.
\]
\end{frame}
%------------------------------------------------
\begin{frame}{Classical Coupled Cluster Ansatz}
In classical coupled cluster theory,
\[
\ket{\Psi_{\mathrm{CC}}}
=
e^T
\ket{\Phi_{\mathrm{HF}}},
\]
where
\[
T = T_1 + T_2 + T_3 + \cdots .
\]
\vspace{0.3cm}
The singles and doubles operators are
\[
T_1
=
\sum_{ia}
t_i^a
a_a^\dagger a_i ,
\]
\[
T_2
=
\frac{1}{4}
\sum_{ijab}
t_{ij}^{ab}
a_a^\dagger a_b^\dagger a_j a_i .
\]
\vspace{0.3cm}
Indices:
\begin{itemize}
\item \(i,j\): occupied spin orbitals,
\item \(a,b\): virtual spin orbitals.
\end{itemize}
\end{frame}
%------------------------------------------------
\begin{frame}{Unitary Coupled Cluster}
The coupled-cluster operator \(e^T\) is not unitary.
\vspace{0.3cm}
In quantum computing, we use the unitary ansatz
\[
\ket{\Psi_{\mathrm{UCC}}}
=
e^{T-T^\dagger}
\ket{\Phi_{\mathrm{HF}}}.
\]
\vspace{0.3cm}
For UCCSD,
\[
T = T_1 + T_2,
\]
so
\[
\ket{\Psi_{\mathrm{UCCSD}}}
=
e^{T_1+T_2-T_1^\dagger-T_2^\dagger}
\ket{\Phi_{\mathrm{HF}}}.
\]
\vspace{0.3cm}
The operator
\[
A = T - T^\dagger
\]
is anti-Hermitian:
\[
A^\dagger = -A.
\]
Therefore \(e^A\) is unitary.
\end{frame}
%------------------------------------------------
\begin{frame}{Single Excitations}
A single excitation promotes one electron:
\[
i \rightarrow a.
\]
The corresponding anti-Hermitian excitation operator is
\[
\tau_i^a
=
a_a^\dagger a_i
-
a_i^\dagger a_a .
\]
The associated unitary is
\[
U_i^a(\theta_i^a)
=
\exp\left[
\theta_i^a
\left(
a_a^\dagger a_i
-
a_i^\dagger a_a
\right)
\right].
\]
\vspace{0.3cm}
This rotates between the reference determinant and singly excited determinants.
\end{frame}
%------------------------------------------------
\begin{frame}{Double Excitations}
A double excitation promotes two electrons:
\[
ij \rightarrow ab.
\]
The anti-Hermitian excitation operator is
\[
\tau_{ij}^{ab}
=
a_a^\dagger a_b^\dagger a_j a_i
-
a_i^\dagger a_j^\dagger a_b a_a .
\]
The associated unitary is
\[
U_{ij}^{ab}(\theta_{ij}^{ab})
=
\exp\left[
\theta_{ij}^{ab}
\tau_{ij}^{ab}
\right].
\]
\vspace{0.3cm}
Doubles are essential for describing electron correlation.
\end{frame}
%------------------------------------------------
\begin{frame}{Trotterized UCCSD Ansatz}
The exact UCCSD ansatz is
\[
U_{\mathrm{UCCSD}}
=
\exp
\left[
\sum_{ia}\theta_i^a \tau_i^a
+
\sum_{ijab}\theta_{ij}^{ab}\tau_{ij}^{ab}
\right].
\]
\vspace{0.3cm}
Because the excitation operators generally do not commute,
\[
[\tau_\mu,\tau_\nu]\neq 0,
\]
one often uses a first-order Trotter approximation:
\[
U_{\mathrm{UCCSD}}
\approx
\prod_{ia}
e^{\theta_i^a\tau_i^a}
\prod_{ijab}
e^{\theta_{ij}^{ab}\tau_{ij}^{ab}}.
\]
\vspace{0.3cm}
The parameters \(\bm{\theta}\) are optimized variationally.
\end{frame}
%------------------------------------------------
\begin{frame}{Quantum Circuit Construction}
The UCCSD circuit is built by applying excitation unitaries to the Hartree--Fock state:
\[
\ket{\Psi(\bm{\theta})}
=
U_{\mathrm{UCCSD}}(\bm{\theta})
\ket{\Phi_{\mathrm{HF}}}.
\]
\vspace{0.3cm}
Workflow:
\begin{enumerate}
\item Prepare \(\ket{\Phi_{\mathrm{HF}}}\) on qubits.
\item Apply all selected single-excitation unitaries.
\item Apply all selected double-excitation unitaries.
\item Measure expectation values of Pauli strings.
\end{enumerate}
\vspace{0.3cm}
The circuit depth can become large because UCCSD contains many excitations.
\end{frame}
%------------------------------------------------
\begin{frame}{Energy Measurement}
After mapping to qubits,
\[
H_q=\sum_\alpha c_\alpha P_\alpha .
\]
For a fixed parameter vector \(\bm{\theta}\), estimate
\[
E(\bm{\theta})
=
\sum_\alpha c_\alpha
\bra{\Psi(\bm{\theta})}
P_\alpha
\ket{\Psi(\bm{\theta})}.
\]
\vspace{0.3cm}
Each Pauli string is measured on the quantum device.
\vspace{0.3cm}
For example,
\[
P_\alpha = Z_0 Z_1 X_2 Y_3.
\]
Measurement requires rotating qubits into the appropriate basis.
\end{frame}
%------------------------------------------------
\begin{frame}{Classical Optimization Loop}
The VQE loop is hybrid:
\[
\bm{\theta}^{(k)}
\longrightarrow
\text{quantum circuit}
\longrightarrow
E(\bm{\theta}^{(k)})
\longrightarrow
\text{classical optimizer}.
\]
\vspace{0.4cm}
The optimizer updates
\[
\bm{\theta}^{(k+1)}
=
\bm{\theta}^{(k)}
+
\Delta \bm{\theta}^{(k)}.
\]
\vspace{0.4cm}
The process continues until
\[
|E^{(k+1)}-E^{(k)}| < \epsilon.
\]
\end{frame}
%------------------------------------------------
\begin{frame}{Algorithm: VQE-UCCSD}
\small
\begin{enumerate}
\item Choose a molecular geometry and one-particle basis.
\item Compute one- and two-body integrals:
\[
h_{pq},\qquad h_{pqrs}.
\]
\item Build the second-quantized Hamiltonian \(H\).
\item Map \(H\) to a qubit Hamiltonian:
\[
H_q=\sum_\alpha c_\alpha P_\alpha.
\]
\item Prepare the Hartree--Fock reference state.
\item Construct the UCCSD ansatz:
\[
U(\bm{\theta})=e^{T_1+T_2-T_1^\dagger-T_2^\dagger}.
\]
\item Measure
\[
E(\bm{\theta})=\bra{\Psi(\bm{\theta})}H_q\ket{\Psi(\bm{\theta})}.
\]
\item Minimize \(E(\bm{\theta})\) classically.
\end{enumerate}
\end{frame}
%------------------------------------------------
\begin{frame}{Variational Principle}
The VQE energy satisfies
\[
E(\bm{\theta})
=
\bra{\Psi(\bm{\theta})}H\ket{\Psi(\bm{\theta})}
\geq
E_0.
\]
\vspace{0.3cm}
Therefore, if the state is normalized,
\[
E_{\mathrm{VQE}}
=
\min_{\bm{\theta}} E(\bm{\theta})
\]
is an upper bound to the exact ground-state energy.
\vspace{0.3cm}
Improving the ansatz lowers the variational energy.
\end{frame}
%------------------------------------------------
\begin{frame}{Gradient Evaluation}
One may use gradient-free optimizers, but gradients can also be computed.
For a parameter \(\theta_\mu\),
\[
\frac{\partial E}{\partial \theta_\mu}
=
\bra{\partial_\mu \Psi}H\ket{\Psi}
+
\bra{\Psi}H\ket{\partial_\mu \Psi}.
\]
\vspace{0.3cm}
For suitable Pauli-generated rotations,
\[
U_\mu(\theta_\mu)
=
e^{-i\theta_\mu P_\mu/2},
\]
one can use the parameter-shift rule:
\[
\frac{\partial E}{\partial \theta_\mu}
=
\frac12
\left[
E\left(\theta_\mu+\frac{\pi}{2}\right)
-
E\left(\theta_\mu-\frac{\pi}{2}\right)
\right].
\]
\end{frame}
%------------------------------------------------
\begin{frame}{Scaling of UCCSD}
Let
\[
N_{\mathrm{occ}} = \text{number of occupied spin orbitals},
\]
\[
N_{\mathrm{virt}} = \text{number of virtual spin orbitals}.
\]
\vspace{0.3cm}
Number of single excitations:
\[
N_{\mathrm{singles}}
\sim
N_{\mathrm{occ}}N_{\mathrm{virt}}.
\]
Number of double excitations:
\[
N_{\mathrm{doubles}}
\sim
N_{\mathrm{occ}}^2N_{\mathrm{virt}}^2.
\]
\vspace{0.3cm}
Thus UCCSD can become expensive for large systems.
\end{frame}
%------------------------------------------------
\begin{frame}{Why UCCSD Is Useful}
Advantages:
\begin{itemize}
\item Physically motivated ansatz.
\item Direct connection to quantum chemistry.
\item Captures dynamical electron correlation.
\item Systematically improvable.
\item Unitary form is natural for quantum circuits.
\end{itemize}
\vspace{0.3cm}
Main limitations:
\begin{itemize}
\item Deep circuits.
\item Many variational parameters.
\item Trotter errors.
\item Measurement overhead.
\item Optimization can be difficult.
\end{itemize}
\end{frame}
%------------------------------------------------
\begin{frame}{Conceptual Summary}
The VQE-UCCSD method combines:
\[
\text{Quantum chemistry}
+
\text{unitary coupled cluster}
+
\text{quantum measurement}
+
\text{classical optimization}.
\]
\vspace{0.4cm}
The central variational problem is
\[
E_{\mathrm{VQE}}
=
\min_{\bm{\theta}}
\bra{\Phi_{\mathrm{HF}}}
U^\dagger_{\mathrm{UCCSD}}(\bm{\theta})
H
U_{\mathrm{UCCSD}}(\bm{\theta})
\ket{\Phi_{\mathrm{HF}}}.
\]
\vspace{0.3cm}
The quantum computer prepares the correlated wave function.
The classical computer updates the cluster amplitudes.
\end{frame}
%------------------------------------------------
\begin{frame}{References}
\footnotesize
\begin{thebibliography}{99}
\bibitem{Peruzzo2014}
A. Peruzzo et al.,
\newblock A variational eigenvalue solver on a photonic quantum processor,
\newblock \emph{Nature Communications} 5, 4213 (2014).
\bibitem{McClean2016}
J. R. McClean et al.,
\newblock The theory of variational hybrid quantum-classical algorithms,
\newblock \emph{New Journal of Physics} 18, 023023 (2016).
\bibitem{Romero2018}
J. Romero et al.,
\newblock Strategies for quantum computing molecular energies using the unitary coupled cluster ansatz,
\newblock \emph{Quantum Science and Technology} 4, 014008 (2018).
\bibitem{Bartlett2007}
R. J. Bartlett and M. Musial,
\newblock Coupled-cluster theory in quantum chemistry,
\newblock \emph{Reviews of Modern Physics} 79, 291 (2007).
\end{thebibliography}
\end{frame}
%------------------------------------------------
\end{document}
